\section{Marco conceptual}
\label{sec:background}

\subsection{Clasificación de imágenes}

La clasificación de imágenes es una tarea de visión por computador donde un modelo
asigna una etiqueta a una imagen de entrada~\cite{lecun2015deep,he2016deep}. En
este proyecto, la entrada es una fotografía de una roca y la salida es una de nueve
clases posibles. El modelo devuelve probabilidades por clase; la aplicación usa la
clase con mayor probabilidad como predicción principal y muestra también
alternativas cercanas.

\subsection{Azure Custom Vision}

Azure Custom Vision es un servicio de Azure AI para entrenar modelos de
clasificación y detección de objetos con un flujo gestionado. Permite cargar
imágenes etiquetadas, entrenar iteraciones, evaluar el desempeño y publicar un
endpoint de predicción~\cite{azurecustomvision}.

Para Litoteca se seleccionó Custom Vision por su facilidad para entrenar un
clasificador multiclase, exponer una API de predicción y separar el entrenamiento
del consumo desde una aplicación web. La aplicación no expone las credenciales al
navegador: una Pages Function en Cloudflare Pages recibe la imagen, consulta la
API de predicción y devuelve una respuesta normalizada al
frontend~\cite{cloudflarefunctions}.

\subsection{Métricas de evaluación}

La evaluación de un clasificador multiclase se apoya en métricas derivadas de la
matriz de confusión. Para cada clase se consideran verdaderos positivos (VP),
falsos positivos (FP) y falsos negativos (FN).

\begin{equation}
  \mathrm{Precision} = \frac{VP}{VP + FP}
  \label{eq:precision}
\end{equation}

\begin{equation}
  \mathrm{Recall} = \frac{VP}{VP + FN}
  \label{eq:recall}
\end{equation}

\begin{equation}
  F1 = 2 \cdot \frac{\mathrm{Precision} \cdot \mathrm{Recall}}{\mathrm{Precision} + \mathrm{Recall}}
  \label{eq:f1}
\end{equation}

La precisión resume qué tan confiables son las predicciones positivas, el recall
mide la cobertura sobre los casos reales de la clase y F1 balancea ambos. En Custom
Vision, la AP (average precision) resume el comportamiento de la curva
precisión-recall por clase.

\subsection{Clases del prototipo}

El prototipo considera nueve clases, seleccionadas por disponibilidad en el dataset
y por su valor didáctico. La Tabla~\ref{tab:clases} resume la taxonomía usada por
la aplicación.

\begin{table}[H]
  \centering
  \caption{Clases consideradas en Litoteca}
  \label{tab:clases}
  \begin{tabular}{@{}llll@{}}
    \toprule
    ID & Nombre común & Nombre técnico & Origen \\
    \midrule
    basalt & Basalto & Basalt & Ígneo \\
    olivine-basalt & Basalto olivínico & Olivine basalt & Ígneo \\
    chert & Sílex & Chert & Sedimentario \\
    clay & Arcilla & Clay & Sedimentario \\
    conglomerate & Conglomerado & Conglomerate & Sedimentario \\
    diatomite & Diatomita & Diatomite & Sedimentario \\
    gypsum & Yeso & Gypsum & Sedimentario \\
    shale & Lutita & Shale (mudstone) & Sedimentario \\
    siliceous-sinter & Geiserita & Siliceous sinter & Sedimentario \\
    \bottomrule
  \end{tabular}
\end{table}
